Label the starting vertex $A=(0,0)$ and the others $B=(2,0)$, $C=(1,\sqrt3)$. The first reflection must be on $BC$. Up to reflection symmetry, the next two sides are $AC,AB$, and the endpoint must be $C$; an endpoint on the last reflecting side would require travelling along that side, contradicting an interior reflection.
Unfold the path by reflecting the endpoint successively across the last, middle, and first reflecting sides. The images of $C$ are $(1,-\sqrt3)$ across $AB$, then $(-2,0)$ across $AC$, then $(4,2\sqrt3)$ across $BC$. Reflections preserve path lengths and make the reflected path straight, so the distance is
\[\boxed{\sqrt{4^2+(2\sqrt3)^2}=2\sqrt7\text{ metres}}.\]
The path is valid: its three bounce points are $(\dfrac{4}{3},\dfrac{2\sqrt3}{3})$, $(\dfrac{1}{2},\dfrac{\sqrt3}{2})$, $(\dfrac{2}{3},0)$, all in side interiors. Folding the straight path back gives the required equal reflection angles.
\begin{asy}
texpreamble("\usepackage{amsmath,amssymb}");
size(240);
defaultpen(fontsize(9pt));
pair A=(0,0),B=(2,0),C=(1,sqrt(3)),Q=(4/3,2sqrt(3)/3),R=(.5,sqrt(3)/2),S=(2/3,0);draw(A--B--C--cycle);draw(A--Q--R--S--C,blue+linewidth(1.2));label("$A$",A,(-1,-1));label("$B$",B,(1,-1));label("$C$",C,(0,1));label("$Q$",Q,(1,0));label("$R$",R,(-1,0));label("$S$",S,(0,-1));
\end{asy}
\begin{asy}
texpreamble("\usepackage{amsmath,amssymb}");
size(240);
defaultpen(fontsize(9pt));
pair A=(0,0),B=(2,0),C=(1,sqrt(3)),U=(1,-sqrt(3)),V=(-2,0),W=(4,2sqrt(3));draw(A--B--C--cycle,gray);draw(C--U--V--W,dashed);draw(A--W,blue+linewidth(1.2));label("$A$",A,(0,-1));label("$C$",C,(0,1));label("$C^{(1)}$",U,(0,-1));label("$C^{(2)}$",V,(-1,0));label("$C^{(3)}(4,2\sqrt3)$",W,(0,1));
\end{asy}
